% part: intuitionistic-logic
% chapter: soundness-completeness
% section: soundness-axd

\documentclass[../../../include/open-logic-section]{subfiles}

\begin{document}

\olfileid{int}{sc}{sax}

\olsection{स्वयंसिद्धकीय \usetoken{P}{derivation} ची निर्दोषता}

\begin{editorial}
  निर्दोषतेची सिद्धता सर्व स्वयंसिद्धके
  अंतःप्रज्ञावादी अर्थाने वैध आहेत, या
  तथ्यावर अवलंबून आहे. हे अजून सिद्ध
  करायचे आहे; उदाहरणार्थ, चिन्हार्थमीमांसेच्या
  प्रकरणात ते करता येईल.
\end{editorial}

\begin{thm}[निर्दोषता]
  \ollabel{thm:soundness}
  $\Gamma \Proves !A$ असेल, तर
  $ \Gamma \Entails !A$.
\end{thm}

\begin{proof}
  $\Gamma \Proves !A$ असेल, तर
  $\Gamma \Entails !A$, हे सिद्ध करू.
  $\Gamma$ पासून~$!A$ च्या
  !!{derivation} मधील !!{formula}s ची
  संख्या~$n$ यावर विगमनाने सिद्धता करतो.
  $!A_1$, \dots, $!A_n = !A$ ही
  $\Gamma$ पासूनची !!a{derivation} असेल,
  तर $\Gamma \Entails !A_n$ हे दाखवू.
  $!A_1$, \dots, $!A_n$ ही
  !!a{derivation} असेल, तर कोणत्याही $k<n$
  साठी $!A_1$, \dots, $!A_k$ हीही
  निष्पत्ती आहे, हे लक्षात घ्या.

  लांबी~$0$ च्या !!{derivation}s नसतात; म्हणून
  $n=0$ साठी दावा रिक्तपणे खरा आहे.
  त्यामुळे लांबी~$<n$ असलेल्या सर्व
  !!{derivation}s साठी दावा खरा आहे, असे
  विगमन-गृहीतक घेऊ. आता~$!A_n$ चे
  समर्थन कशावर आहे त्यानुसार प्रकार पाहू.

  \begin{enumerate}
  \item $!A_n$ हे स्वयंसिद्धक आहे.
    सर्व स्वयंसिद्धके वैध असल्याने कोणत्याही~$\Gamma$
    साठी $\Gamma \Entails
    !A_n$.
  \item $!A_n \in \Gamma$. मग कोणतेही
    $\mModel{M}$ आणि $w$ घ्या.
    $\mSat{M}{\Gamma}[w]$ असेल, तर उघडच
    $\mSat{M}{!A_n}[w]$; म्हणजे
    $\Gamma \Entails !A$.
  \item $!A_n$ हे $!A_i$ आणि
    $!A_j \ident !A_i \lif
    !A_n$ यांपासून \MP{} ने निष्पन्न होते.
    $!A_1$, \dots, $!A_i$ आणि
    $!A_1$, \dots, $!A_j$ या
    $\Gamma$ पासूनच्या !!{derivation}s आहेत.
    म्हणून विगमन-गृहीतकानुसार
    $\Gamma \Entails !A_i$ आणि
    $\Gamma \Entails !A_i \lif !A_n$.

    $\mSat{M}{\Gamma}[w]$ असे समजा.
    $\mSat{M}{\Gamma}[w]$ आणि
    $\Gamma \Entails !A_i \lif !A_n$
    असल्यामुळे $\mSat{M}{!A_i \lif
      !A_n}[w]$. व्याख्येनुसार याचा अर्थ
    असा की $Rww'$ असलेल्या प्रत्येक $w'$
    साठी $\mSat{M}{!A_i}[w']$ असेल, तर
    $\mSat{M}{!A_n}[w']$. $R$ परावर्ती
    असल्याने $w$ हेही $Rww'$ पूर्ण करणाऱ्या
    $w'$ पैकी एक आहे. म्हणजे
    $\mSat{M}{!A_i}[w]$ असेल, तर
    $\mSat{M}{!A_n}[w]$.
    $\Gamma \Entails !A_i$ असल्याने
    $\mSat{M}{!A_i}[w]$. म्हणून
    अपेक्षेप्रमाणे $\mSat{M}{!A_n}[w]$.
  \end{enumerate}
\end{proof}

\end{document}
